\section{Artefato e caso empírico}

Este capítulo documentará o sistema RAG e o corpus de discursos da 56\textordfeminine{} Legislatura do Senado Federal empregados como caso empírico. A prova de conceito já implementada constitui o ponto de partida técnico; a versão efetivamente avaliada na dissertação deverá ser descrita a partir de um congelamento reproduzível de dados, código, modelos, parâmetros e configurações.

\subsection{Caracterização do corpus legislativo}

% TODO: registrar origem, cobertura temporal, critérios de inclusão e exclusão,
% volume final, campos, duplicatas, ausências e hashes do congelamento utilizado.

\subsection{Arquitetura e componentes do sistema RAG}

% TODO: adaptar a caracterização e a proposta do artigo-base, atualizando apenas
% configurações confirmadas na execução final da dissertação.

\subsection{Preparação, segmentação e enriquecimento documental}

% TODO: descrever chunking, sobreposição, metadados, embeddings e ingestão.

\subsection{Configuração experimental e reprodutibilidade}

% TODO: registrar versões, modelos, prompts, top-k, reranqueamento, temperatura,
% datas de execução, sementes quando aplicáveis e localização dos artefatos.

\subsection{Evolução em relação à prova de conceito}

% TODO: separar claramente o que foi herdado do artigo, o que foi modificado e
% o que foi criado especificamente para a dissertação.
